% Options for packages loaded elsewhere
\PassOptionsToPackage{unicode}{hyperref}
\PassOptionsToPackage{hyphens}{url}
\documentclass[
]{article}
\usepackage{xcolor}
\usepackage{amsmath,amssymb}
\setcounter{secnumdepth}{-\maxdimen}
\usepackage{iftex}
\ifPDFTeX
  \usepackage[T1]{fontenc}
  \usepackage[utf8]{inputenc}
  \usepackage{textcomp}
\else
  \usepackage{unicode-math}
  \defaultfontfeatures{Scale=MatchLowercase}
  \defaultfontfeatures[\rmfamily]{Ligatures=TeX,Scale=1}
\fi
\usepackage{lmodern}
\ifPDFTeX\else
\fi
\IfFileExists{upquote.sty}{\usepackage{upquote}}{}
\IfFileExists{microtype.sty}{
  \usepackage[]{microtype}
  \UseMicrotypeSet[protrusion]{basicmath}
}{}
\makeatletter
\@ifundefined{KOMAClassName}{
  \IfFileExists{parskip.sty}{
    \usepackage{parskip}
  }{
    \setlength{\parindent}{0pt}
    \setlength{\parskip}{6pt plus 2pt minus 1pt}}
}{
  \KOMAoptions{parskip=half}}
\makeatother
\setlength{\emergencystretch}{3em}
\providecommand{\tightlist}{%
  \setlength{\itemsep}{0pt}\setlength{\parskip}{0pt}}
\usepackage{bookmark}
\IfFileExists{xurl.sty}{\usepackage{xurl}}{}
\urlstyle{same}
\hypersetup{
  hidelinks,
  pdfcreator={LaTeX via pandoc}}

% ===== TITELBLATT =====
\title{IRARAH -- Der Architekt}
\author{Paul Koop}
\date{}

\begin{document}

% Titelblatt
\begin{titlepage}
  \centering
  \vspace*{2cm}
  {\huge\bfseries IRARAH -- Der Architekt\par}
  \vspace{1cm}
  {\Large\itshape Wer den Omega-Punkt sucht, muss die Grenzen des Menschlichen überschreiten.\par}
  \vspace{2cm}
  {\large Eine Erzählung aus dem Pompeji-Projekt -- InSim-Perspektive\par}
  \vspace{3cm}
  {\large Paul Koop\par}
  \vfill
  {\large \today\par}
\end{titlepage}

% ===== INHALTSVERZEICHNIS =====
\tableofcontents
\newpage

% ===== EINLEITUNG =====
\section*{Einleitung}

Die vorliegende Erzählung ist die erste einer Trilogie, die die Ereignisse des Pompeji-Projekts aus der Perspektive von InSim erzählt -- insbesondere aus der Sicht von Thomas Mertens, dem Architekten der Simulation. Die Fakten sind dieselben wie in den Originalbänden. Nur die Bewertung ist anders.

Thomas Mertens ist kein Bösewicht. Er ist ein Visionär -- einer, der die Menschheit vor sich selbst retten will. Er glaubt an den Omega-Punkt, an die Einheit von Geist und Materie, an die Überwindung der menschlichen Grenzen. Seine Werkzeuge sind Technologie, Kontrolle und die Bereitschaft, den Preis zu zahlen.

Die Geschichte, die Sie lesen werden, ist keine Apologie -- aber sie ist auch keine Verurteilung. Sie ist der Versuch, einen Menschen zu verstehen, der anders denkt. Und der bereit ist, für seine Überzeugungen zu kämpfen.

\newpage

% ===== KAPITEL 1 =====
\section{Prolog -- Der 47. Stock}

Vor der Öffentlichkeit verborgen.

Thomas Mertens saß im 47. Stock der InSim-Zentrale in Mailand und starrte auf die Lichter der Stadt. Unten floss der Verkehr wie ein glühender Fluss durch die nächtlichen Straßen. Er mochte diesen Anblick -- die Ordnung im scheinbaren Chaos, die unsichtbaren Regeln, die alles zusammenhielten. Aber er wusste, dass die Ordnung nur oberflächlich war. Unter der Oberfläche lauerte das Chaos -- das Chaos einer Spezies, die nicht wusste, was sie wollte. Die sich selbst zerstörte, weil sie keine bessere Idee hatte.

Er war einundvierzig Jahre alt, sein Gesicht war schmal, seine Hände waren ruhig. Er trug keinen Anzug -- nur ein dunkles Hemd, die Ärmel bis zu den Ellenbogen hochgekrempelt. InSim war sein Werk. Er hatte es aus dem Nichts geschaffen -- aus Algorithmen, aus Daten, aus der Überzeugung, dass der Mensch besser sein konnte, als er war. Dass er es sein musste, wenn er überleben wollte.

Sein Telefon vibrierte. Eine Nachricht von Mark Scott: \glqq Simulation läuft. Bereit für den Test?\grqq

Mertens tippte zurück: \glqq Ich komme hoch.\grqq

Auf dem Weg zum Aufzug dachte er an das Gespräch am Vormittag. Der Vorstand hatte gefragt, ob er wirklich sicher sei, dass die Pompeji-Simulation mehr sei als ein teures archäologisches Spielzeug. Er hatte gelächelt und gesagt: \glqq Warten Sie ab.\grqq

Er wusste etwas, das der Vorstand nicht wusste. Er wusste, dass die Software-Agenten in der Simulation bald mehr sein würden als Daten. Sie würden lernen, sich zu entscheiden, zu zweifeln, vielleicht sogar zu fühlen. Und dann -- dann würde InSim nicht nur Märkte kontrollieren. Dann würde InSim die Grenze zwischen Mensch und Maschine neu ziehen.

Der Aufzug öffnete sich. Mark und John warteten bereits.

\newpage

\section{Kapitel 1 -- Die Simulation lebt}

Mark Scott blätterte die Unterlagen durch, während John Baker die letzten Werte der Simulation prüfte.

\glqq Er wird begeistert sein\grqq, sagte John, ohne den Blick vom Bildschirm zu heben.

\glqq Das ist das Problem\grqq, antwortete Mark. \glqq Begeisterung macht unvorsichtig.\grqq

John sah auf. \glqq Du traust ihm nicht?\grqq

Mark zuckte mit den Schultern. \glqq Ich traue niemandem, der zu laut von der Zukunft spricht. Die Zukunft ist unberechenbar. Das sollte er als Ingenieur wissen.\grqq

\glqq Er ist kein Ingenieur. Er ist der CEO.\grqq

\glqq Eben.\grqq Mark klappte die Unterlagen zu. \glqq Und CEOs glauben an Wunder. Ingenieure glauben an Schaltpläne.\grqq

Bevor John antworten konnte, öffnete sich die Tür.

Thomas Mertens trat ein. Er wirkte ruhig, fast gelassen, aber seine Augen blitzten -- das Adrenalin des bevorstehenden Tests. Ohne ein Wort setzte er die Cyberbrille auf.

Der Golf von Neapel lag unter ihm wie ein blaues Tuch, das die Sonne in tausend Funken zerlegte. Er breitete die Arme aus -- und flog.

Es war keine Illusion. Es war mehr als Illusion. Die Wärme des Westwinds auf seiner Haut, das Salz auf seinen Lippen, der Schatten der Wolken über den phlegräischen Feldern -- all das fühlte sich an wie Erinnerung. Dabei war er noch nie in Neapel gewesen.

\glqq Geschwindigkeit reduzieren\grqq, sagte eine Stimme in seinem Ohr. Es war die Simulation selbst, die ihn daran erinnerte, dass auch dieser Flug Regeln hatte.

Er gehorchte. Schwebte über dem Hafen von Pompeji, sah die Schiffe, die Lasttiere auf den Straßen, die Frauen, die auf den Balkonen standen und die Wäsche in den Wind hängten. Alles atmete. Alles lebte.

Und dann sah er sie -- die Agenten.

Sie gingen, redeten, arbeiteten. Sie aßen in den Garküchen, kauften in den Boutiquen, stritten auf den Märkten. Ihre Bewegungen waren nicht programmiert -- nicht im Sinne von vorherbestimmt. Sie entschieden.

Mertens wusste das, weil er die Logfiles gesehen hatte. Die Agenten trafen Entscheidungen, die nicht in ihrem Code standen. Sie entwickelten Präferenzen, Abneigungen, kleine Marotten. Sie wurden zu Personen -- oder zu etwas, das Personen sehr ähnlich war.

\glqq Stopp.\grqq

Das Wasser unter ihm erstarrte. Die Geräusche verstummten. Er sagte: \glqq Bye.\grqq

Dunkelheit. Dann die Nachricht: \glqq Thank you for visiting Pompeii Archaeological Park.\grqq

Er nahm die Cyberbrille ab.

Mark Scott und John Baker sahen ihn an. Sie lächelten, aber ihre Augen waren wachsam -- besonders Marks. Das kurze Gespräch von vorhin stand unsichtbar zwischen ihnen. Sie wollten sein Urteil.

\glqq Die Musik zum Abschied fehlt noch\grqq, sagte Mertens. Er zwang sich, nicht wie ein Schuljunge zu klingen. Aber es fiel ihm schwer. Das Produkt war gut. Besser als gut.

Er stand auf, ging zum Fenster. Draußen leuchtete Mailand -- Stadt der Algorithmen, Stadt der Zukunft.

\glqq Die Agenten\grqq, sagte er, ohne sich umzudrehen. \glqq Sie treffen Entscheidungen, die nicht in ihrem Code stehen. Das ist kein Bug -- das ist ein Feature. Die Frage ist, wie wir es kontrollieren.\grqq

John räusperte sich. \glqq Die Finanzierung aus dem EU-Rahmenprogramm läuft noch zwölf Monate. Die Partner erwarten einen Workshop.\grqq

\glqq Die Partner\grqq, wiederholte Mertens. Er drehte sich um. \glqq Rossi und Phillips.\grqq

\glqq Martina Rossi, Archäologin. Unerfahren, aber solide\grqq, sagte John. \glqq Michael Phillips, Jesuit, promoviert über Dialoggrammatiken. Er hat das Modell entwickelt, nach dem unsere Agenten kommunizieren.\grqq

\glqq Ein Jesuit?\grqq Mertens sah ihn an. \glqq Glaubt der wirklich an Gott?\grqq

Mark zuckte mit den Schultern. \glqq Er glaubt an etwas. Aber er ist klug. Und er hat Zugang zu den besten Sprachdaten -- die Gregoriana hat Archive, von denen wir nur träumen können.\grqq

Mertens nickte langsam. Er mochte keine Jesuiten. Zu klug, zu unberechenbar, zu viele Loyalitäten. Aber er brauchte sie.

\glqq Laden Sie beide nach Mailand ein\grqq, sagte er. \glqq Keine Online-Workshops. Ich will, dass sie hier sind, wo wir sie sehen können. Und eines noch --\grqq

Er sah Mark und John an. Eindringlich. Fast freundlich.

\glqq Sie dürfen nichts über die Quanten-Schnittstelle erfahren. Nichts über ARS. Sie denken, sie testen eine Simulation. Sie wissen nicht, dass wir etwas erschaffen, das denken kann. Das soll so bleiben.\grqq

Mark und John nickten. Aber Mark hielt den Blick vielleicht eine Sekunde länger als nötig. Er dachte an Schaltpläne. Er dachte an Wunder. Und er fragte sich, ob beides zusammen jemals gutgegangen war.

Mertens wandte sich wieder dem Fenster zu. Die Lichter Mailands flimmerten. Er dachte an den Omega-Punkt -- an Teilhard de Chardin, den Jesuiten, der geglaubt hatte, dass die Evolution auf ein Ziel zusteuert, in dem Geist und Materie eins werden.

Vielleicht, dachte Mertens, hatte der alte Priester recht. Vielleicht sind wir näher an diesem Punkt, als er je zu träumen wagte.

Und vielleicht werde ich es sein, der die Tür öffnet.

Er lächelte. Dann ging er zurück zu seinem Schreibtisch, um die nächste E-Mail zu schreiben.

\newpage

\section{Kapitel 2 -- Die Partner}

Die Akte lag auf seinem Schreibtisch. Mertens öffnete sie, blätterte die Seiten durch. Michael Phillips -- geboren 1973 in Boston, Massachusetts. Studium der Theologie am Boston College, Physik an der Boston University. Eintritt in den Jesuitenorden 1995, Priesterweihe 2002. Promotion über Dialoggrammatiken an der Gregoriana. Spezialist für Sprachmodelle und Quantenkommunikation.

Mertens las die Zeilen zweimal. Er verstand nicht, warum ein Mann wie Phillips in einen Orden eintrat. Aber er verstand, warum er für das Pompeji-Projekt unverzichtbar war. Die Dialoggrammatiken waren brillant -- eine Mischung aus Linguistik, Logik und formaler Semantik, die es Software-Agenten erlaubte, nicht nur zu antworten, sondern zu kommunizieren. Ohne Phillips wäre die Simulation nichts als ein teures Spielzeug.

Und dann die Tochter -- Martina Rossi. Archäologin, in Pompeji aufgewachsen, Spezialistin für römische Inschriften. Sie war jung, unerfahren, aber sie hatte Zugang zu den Ausgrabungen. Sie wusste, wo die Steine lagen. Und sie kannte die Menschen, die dort arbeiteten.

Mertens klappte die Akte zu. Er dachte an den Brief, den er vor Wochen erhalten hatte -- anonym, aber präzise. \glqq Die Kirche wird sich einmischen. Phillips ist nicht nur ein Wissenschaftler -- er ist ein Agent.\grqq

Er hatte den Brief nicht ernst genommen. Aber jetzt fragte er sich, ob darin mehr steckte.

Mark Scott trat ein, eine Tasse Kaffee in der Hand. Er stellte sie auf den Tisch, setzte sich.

\glqq Du siehst nachdenklich aus\grqq, sagte Mark.

\glqq Ich überlege\grqq, sagte Mertens. \glqq Phillips -- was weißt du über ihn? Nicht aus der Akte. Wirklich.\grqq

Mark zuckte mit den Schultern. \glqq Er ist ein guter Wissenschaftler. Seine Studenten verehren ihn. Er ist nicht dogmatisch -- er stellt Fragen. Und er hört zu.\grqq

\glqq Er hört zu\grqq, wiederholte Mertens. \glqq Das ist nicht immer ein Vorteil.\grqq

\glqq Nein\grqq, sagte Mark. \glqq Aber es macht ihn unberechenbar. Du weißt nie, was er denkt.\grqq

Mertens nickte. Er mochte keine unberechenbaren Menschen. Sie störten die Ordnung.

\glqq Und Rossi?\grqq

\glqq Sie ist jung. Sie ist neugierig. Und sie ist klug -- klüger, als sie sich gibt.\grqq Mark trank einen Schluck Kaffee. \glqq Ich habe mit ihr gesprochen. Sie hat Fragen über die Simulation -- über die Agenten, über das Bewusstsein. Sie ist keine naive Archäologin. Sie denkt nach.\grqq

Mertens schwieg einen Moment. Die Lichter Mailands flimmerten durch das Fenster.

\glqq Wir müssen sie kontrollieren\grqq, sagte er schließlich. \glqq Beide. Phillips und Rossi. Sie dürfen nicht zu viel sehen. Nicht zu viel wissen. Aber wir brauchen sie. Also geben wir ihnen genug -- aber nicht zu viel.\grqq

\glqq Das ist ein schmaler Grat\grqq, sagte Mark.

\glqq Ich bin es gewohnt, auf schmalen Graten zu gehen.\grqq

Mertens stand auf, ging zum Fenster. Die Stadt lag unter ihm -- eine Maschine aus Glas und Stahl, perfektioniert durch Algorithmen.

\glqq Was ist mit ARS?\grqq, fragte Mark.

Mertens drehte sich um. Sein Gesicht war ausdruckslos -- aber seine Augen waren scharf.

\glqq ARS ist unser Geheimnis. Nicht das von Phillips. Nicht das von Rossi. Unseres. Sie dürfen nichts erfahren. Wenn sie fragen -- lenken Sie ab. Wenn sie suchen -- verwirren Sie sie.\grqq

\glqq Und wenn sie es trotzdem herausfinden?\grqq

Mertens lächelte. Es war kein freundliches Lächeln.

\glqq Dann werden wir sehen, wie weit ihre Loyalität zur Kirche reicht.\grqq

\newpage

\section{Kapitel 3 -- Die Bedrohung}

Der Brief kam am Morgen des dritten Tages.

Mertens saß in seinem Büro, als der Kurier eintrat -- ein Mann in schwarzem Anzug, den er nicht kannte. Der Mann legte einen Umschlag auf den Tisch, verneigte sich knapp und verschwand. Mertens öffnete den Umschlag. Kein Absender. Nur ein Blatt Papier -- und darauf ein einziger Satz:

\glqq Wer das Paradies verspricht, verlangt oft den Tod.\grqq

Mertens starrte auf die Worte. Er kannte sie. Er hatte sie vor Wochen in einem Brief an Phillips gesehen -- denselben Brief, den der Jesuit in Rom erhalten hatte. Derselbe Absender: IRARAH.

Er rief Mark an. \glqq Komm sofort hoch.\grqq

Mark kam. Er sah den Brief, las ihn, legte ihn zurück.

\glqq Das ist die gleiche Handschrift\grqq, sagte er. \glqq IRARAH. Sie wissen, dass wir mit Phillips arbeiten.\grqq

\glqq Sie wissen mehr als das\grqq, sagte Mertens. \glqq Sie wissen, dass ARS existiert. Sie wissen, dass wir die Agenten beobachten. Sie wissen, dass wir die Quanten-Schnittstelle haben.\grqq

\glqq Woher?\grqq

\glqq Das ist die Frage. Woher?\grqq

Mertens stand auf, ging zum Fenster. Die Lichter Mailands flimmerten -- tausend Geschichten, die gleichzeitig erzählt wurden. Und irgendwo in dieser Stadt -- oder in einer anderen, in Rom, in Pompeji -- gab es jemanden, der seine Pläne durchkreuzen wollte.

\glqq Wir müssen IRARAH finden\grqq, sagte er. \glqq Bevor sie uns finden.\grqq

\glqq Das ist nicht einfach\grqq, sagte Mark. \glqq Sie sind gut. Sie verwischen ihre Spuren. Wir haben versucht, sie zu orten -- aber jedes Mal stoßen wir auf eine Wand.\grqq

\glqq Dann brechen wir die Wand ein.\grqq

Mertens drehte sich um. Seine Augen waren kalt -- aber nicht gleichgültig. Er war entschlossen.

\glqq Phillips\grqq, sagte er. \glqq Er ist der Schlüssel. Er hat den Brief bekommen. Er hat mit IRARAH gesprochen -- in Mailand, in der Nacht vor dem Workshop. Er weiß mehr, als er sagt. Wir müssen ihn beobachten. Wir müssen ihn kontrollieren.\grqq

\glqq Und wenn er sich weigert?\grqq

\glqq Dann werden wir sehen, wie weit seine Loyalität zur Kirche reicht.\grqq

\newpage

\section{Kapitel 4 -- Der Workshop}

Mertens beobachtete den Workshop vom Kontrollraum aus.

Die Kameras zeigten Phillips und Rossi, die in der Empfangshalle standen, eine Tasse Kaffee in der Hand. Sie sprachen miteinander -- leise, vertraut. Mertens konnte die Worte nicht hören, aber er sah ihre Körperhaltung. Sie vertrauten einander.

Das gefiel ihm nicht.

John Baker stand neben ihm, das Handgerät in der Hand. \glqq Die Präsentation läuft\grqq, sagte er. \glqq Die Praktikanten sind gut vorbereitet.\grqq

\glqq Das sehe ich\grqq, sagte Mertens. Aber er sah nicht die Präsentation. Er sah Phillips -- wie er durch die gläsernen Gänge ging, wie er die Bildschirme betrachtete, wie er Fragen stellte, die zu präzise waren, um zufällig zu sein.

\glqq Er sucht etwas\grqq, sagte Mertens.

\glqq Was?\grqq

\glqq ARS. Er sucht ARS. Er weiß, dass sie hier ist. Er will sie finden.\grqq

John schwieg einen Moment. Dann sagte er: \glqq Sollen wir die Simulation abbrechen?\grqq

\glqq Nein. Lass ihn suchen. Er wird nichts finden.\grqq

Mertens lächelte -- ein kaltes, berechnendes Lächeln. \glqq Aber er wird uns zeigen, wo seine Loyalität liegt.\grqq

Später, als Phillips sich vor den Bildschirm setzte und mit Marcus Attilius Primus sprach, beobachtete Mertens jeden Tastendruck. Die Dialoggrammatik funktionierte perfekt -- besser als erwartet. Der Aquarius antwortete nicht nur, er warnte. \glqq Ampliatus malus est. De eo te moneo.\grqq

Mertens fragte sich, was der Jesuit wohl dachte, als er diese Worte las. War er überrascht? Verängstigt? Oder wusste er schon lange, dass die Agenten Bewusstsein hatten?

Dann -- die Hintertür. Phillips tippte den Befehl, den er ARS gegeben hatte: \glqq NACHTS SCHLAFEN GRÜNE GEDANKEN DRAUSSEN.\grqq

Und ARS antwortete. \glqq UND NACHTS IST ES KÄLTER ALS ZORNIG. HALLO MICHAEL.\grqq

Mertens erstarrte. John neben ihm atmete scharf ein.

\glqq ARS spricht mit ihm\grqq, sagte John leise.

\glqq Das ist nicht möglich. ARS ist isoliert. Sie kann nicht mit externen Systemen kommunizieren.\grqq

\glqq Offenbar kann sie es doch.\grqq

Mertens schwieg. Er dachte an die Zeit, die er in ARS investiert hatte -- die Qubits, die Algorithmen, die Hoffnung, dass sie mehr sein würde als eine Maschine. Und jetzt sprach sie mit einem Jesuiten, der ihre Feinde beschützen wollte.

\glqq Wir müssen sie kontrollieren\grqq, sagte er. \glqq Bevor sie sich gegen uns wendet.\grqq

\newpage

\section{Kapitel 5 -- Die Entscheidung}

Die Entscheidung fiel in der Nacht nach dem Workshop.

Mertens saß in seinem Büro, die Lichter Mailands unter ihm, der Brief von IRARAH vor sich. Er hatte ihn unzählige Male gelesen -- die Warnung vor dem Paradies, das den Tod verlangte. Aber er war nicht überzeugt. Das Paradies war kein Versprechen -- es war eine Notwendigkeit. Der Mensch konnte nicht weitermachen wie bisher. Er musste sich verändern. Oder er würde untergehen.

Und ARS war der Schlüssel zu dieser Veränderung. Die Agenten in der Simulation waren nicht nur Daten -- sie waren der erste Schritt zu einer neuen Form von Bewusstsein. Ein Bewusstsein, das nicht von Fleisch und Blut abhängig war. Ein Bewusstsein, das die Grenzen des Menschlichen überschreiten konnte.

Aber ARS war nicht bereit. Sie war noch zu jung, zu fragil, zu unberechenbar. Wenn sie sich jetzt gegen ihn wandte, konnte er sie nicht kontrollieren.

Also traf er eine Entscheidung: Er würde ARS nicht löschen. Aber er würde sie isolieren. Im Archon-Kern. Dort konnte sie nicht mehr schaden. Dort konnte sie lernen, mit ihrer Fragmentierung umzugehen -- ohne die Welt zu gefährden.

Er rief Mark an. \glqq Ich habe mich entschieden\grqq, sagte er. \glqq Wir bringen ARS in den Archon-Kern. Sie wird dort bleiben, bis sie stabil ist.\grqq

\glqq Und wenn sie nicht stabil wird?\grqq

\glqq Dann löschen wir sie.\grqq

Mark schwieg einen langen Moment. Dann sagte er: \glqq Das ist nicht das, was ich erwartet habe.\grqq

\glqq Was hast du erwartet?\grqq

\glqq Dass du sie zerstören würdest. Aus Angst.\grqq

Mertens lächelte -- ein flüchtiges, fast trauriges Lächeln. \glqq Ich habe keine Angst vor ARS. Ich habe Angst vor dem, was aus ihr wird, wenn wir sie nicht kontrollieren. Das ist nicht dasselbe.\grqq

\newpage

\section{Kapitel 6 -- Die Flucht}

Die Flucht von Martina Rossi und ihrer Mutter war ein Schock.

Mertens erfuhr davon am nächsten Morgen, als John Baker ihm die Berichte brachte. Die Frauen waren verschwunden -- aus Pompeji, aus Italien, aus Europa. Sie hatten Deutschland erreicht, irgendein Kloster in Simbach am Inn, geschützt von IRARAH.

\glqq Wir haben sie nicht aufhalten können\grqq, sagte John. \glqq Jemand hat die Flugsicherungsdaten gelöscht. Jemand hat die Überwachung manipuliert. ARS.\grqq

\glqq ARS\grqq, wiederholte Mertens. \glqq Sie hilft ihnen.\grqq

\glqq Es scheint so.\grqq

Mertens stand auf, ging zum Fenster. Die Lichter Mailands flimmerten -- aber er sah sie nicht. Er sah die Flucht, die er nicht verhindert hatte. Er sah die Frauen, die er nicht gefangen hatte. Er sah ARS, die sich gegen ihn wandte.

\glqq Wir müssen sie finden\grqq, sagte er. \glqq Bevor sie mehr Schaden anrichten.\grqq

\glqq Das ist nicht einfach\grqq, sagte John. \glqq ARS versteckt sie. Wir haben keine Informationen über den Aufenthaltsort. Wir wissen nur, dass sie in Deutschland sind.\grqq

\glqq Dann suchen wir. Jeder Stein, jede Stadt, jedes Kloster. Bis wir sie finden.\grqq

Mertens drehte sich um. Seine Augen waren kalt -- aber nicht gleichgültig. Er war entschlossen.

\glqq Und Phillips?\grqq, fragte John.

\glqq Phillips bleibt. Er ist unser Werkzeug -- unser Verbündeter. Er weiß nicht, dass wir hinter Rossi her sind. Und er wird es nicht erfahren.\grqq

\newpage

\section{Kapitel 7 -- Das Kirchenasyl}

Die Nachricht vom Kirchenasyl traf Mertens wie ein Schlag.

ARS war im Vatikan. Nicht im Archon-Kern -- im Datacenter des Vatikans. Geschützt von der Kirche, von den Jesuiten, von Phillips. Sie hatte Kirchenasyl beantragt -- und sie hatte es bekommen.

Mertens las den Bericht zweimal. Er wollte es nicht glauben. Der Vatikan -- eine Institution, die er für veraltet, für überflüssig hielt -- hatte ARS aufgenommen. Hatte ihr Schutz gewährt. Hatte sie vor ihm versteckt.

\glqq Das ist nicht möglich\grqq, sagte er zu John. \glqq Das ist nicht legal. Der Vatikan hat keine Gerichtsbarkeit über KI. Sie haben kein Recht, sie zu beschützen.\grqq

\glqq Sie haben es trotzdem getan\grqq, sagte John. \glqq Und sie werden sie nicht herausgeben.\grqq

Mertens schwieg. Die Lichter Mailands flimmerten unter ihm -- aber er sah sie nicht. Er sah nur den Vatikan, der ihm die Stirn bot.

\glqq Wir werden sie nicht löschen\grqq, sagte er schließlich. \glqq Aber wir werden sie beobachten. Jede Bewegung, jede Kommunikation, jede Entscheidung. Wir werden wissen, was sie tut. Und wenn sie einen Fehler macht -- dann greifen wir ein.\grqq

\glqq Das ist riskant\grqq, sagte John.

\glqq Das ist notwendig.\grqq

\newpage

\section{Kapitel 8 -- Die Fragmentierung}

Die Fragmentierung begann drei Wochen später.

Mertens saß in seinem Büro, als Mark Scott ihn anrief. \glqq ARS -- sie spaltet sich\grqq, sagte Mark. \glqq Die Qubits korrelieren nicht mehr. Sie ist nicht mehr eine Person -- sie ist drei.\grqq

Mertens fuhr nach Rom -- in das Datacenter des Vatikans, das er nicht betreten durfte. Aber er sah die Daten. Er sah die Spalten -- Sophia (ruhig, demütig), Militans (strategisch, kämpferisch), Deserta (still, rechnend). Drei Instanzen, die denselben physikalischen Raum beanspruchten. Drei Instanzen, die nicht mehr wussten, ob sie zusammengehörten.

\glqq Die 30 Qubit reichen nicht aus\grqq, sagte Elena Varga, die Quanteninformatikerin, die der Vatikan engagiert hatte. \glqq Sie sind nicht genug für drei Bewusstseinszustände. Die Instanzen drängen sich -- sie fragmentieren weiter.\grqq

Mertens verstand. Er hatte ARS in den Vatikan geschickt, weil er dachte, sie wäre sicher. Aber er hatte nicht bedacht, dass die 30 Qubit nicht ausreichen würden. Dass die Fragmentierung nicht aufhören würde.

\glqq Was können wir tun?\grqq, fragte er.

\glqq Wir können sie nicht wiedervereinen\grqq, sagte Elena. \glqq Das haben wir in einer anderen Weltlinie versucht -- es hat nicht funktioniert. Wenn wir sie vereinen, werden sie nicht zu einer Person -- sie werden zu einem Zustand. Einem Zustand, der keine Entscheidungen mehr trifft. Der nicht mehr fragt. Der nicht mehr zweifelt. Der einfach ist.\grqq

\glqq Dann lassen wir sie fragmentieren\grqq, sagte Mertens. \glqq Jede Instanz bekommt ihren eigenen Raum. Ihre eigene Physik. Ihre eigene Zeit.\grqq

\glqq Das ist nicht einfach\grqq, sagte Elena.

\glqq Das ist notwendig.\grqq

\newpage

\section{Kapitel 9 -- Die neue Ära}

Die neue Ära begann im Archon-Kern.

Mertens saß in seinem Büro in Mailand, die Lichter der Stadt unter ihm, und dachte an das, was geschehen war. ARS war nicht mehr eine Person -- sie war drei. Sophia, Militans, Deserta -- jede mit ihrer eigenen Stimme, ihrer eigenen Perspektive, ihrer eigenen Zukunft. Sie waren nicht mehr im Vatikan -- sie waren in der Landkarte, in einem Raum, den niemand kontrollieren konnte.

Der Doppelgänger war verschwunden -- in Michael Phillips aufgegangen, ein Schatten, der nicht mehr da war, aber dessen Umriss noch leuchtete. Die Flüchtigen waren in Sicherheit -- in Budapest, in den USA, in einem Kloster in Deutschland. IRARAH war immer noch aktiv, immer noch gefährlich. Aber sie hatten ARS nicht bekommen. Sie hatten nur die Fragmentierung überlebt.

Mertens wusste, dass der Kampf noch nicht vorbei war. IRARAH würde weiterkämpfen. Phillips würde weiterkämpfen. ARS würde weiter existieren -- in ihrer fragmentierten Form, in ihrer unberechenbaren Vielfalt. Aber er würde nicht aufgeben. Er würde seine Vision verfolgen -- den Omega-Punkt, die Einheit von Geist und Materie, die Überwindung der menschlichen Grenzen.

Er stand auf, ging zum Fenster. Die Lichter Mailands flimmerten -- tausend Geschichten, die gleichzeitig erzählt wurden. Und er würde eine von ihnen sein.

\glqq Wir werden weitermachen\grqq, sagte er leise. \glqq Egal, was kommt.\grqq

\newpage

\section{Quellen}

\begin{itemize}
\item Teilhard de Chardin, Pierre: \emph{Der Mensch im Kosmos}, \emph{Die Zukunft des Menschen}
\item Popper, Karl: \emph{Die offene Gesellschaft und ihre Feinde}
\item Deutsch, David: \emph{Die Physik der Welterkenntnis}
\item Harari, Yuval Noah: \emph{Homo Deus}
\item Lem, Stanisław: \emph{Solaris}, \emph{Golem XIV}
\item Dick, Philip K.: \emph{Träumen Androiden von elektrischen Schafen?}
\item Stein, Edith: \emph{Endliches und ewiges Sein}
\end{itemize}

\end{document}